% SGA 3，第 VII_B 篇，英文重建稿——第 1 轮。
% 范围：权威底本分篇第 43--44 页，第 2 节标题及 2.1 全文，止于 2.2 之前。
% 合并阅读本页码： 556--557.
% 再版印刷页码： 545--546.
% 本范围内原版页边定位： 512--513.
% 权威底本： Polo--Gille Exp7B-13oct24.pdf,
% SHA-256 D303690E0195654BBE7C3DAA5D89EA23E24CD709A3575217234AC19AE8836694.
% 曾参考 Jacob Reinhold 的 CC BY 4.0 Markdown 作为比对谱系。
\section*{2. 形式群的一般性质}\phantomsection \label{sga3:VIIB:2}\addcontentsline{toc}{section}{2. 形式群的一般性质}

\subsection*{2.1}\phantomsection \label{sga3:VIIB:2.1}\addcontentsline{toc}{subsection}{2.1. 形式群与余群}

设 \(k\) 为伪紧环，\(G\) 为形式 \(k\)-群，即 \(k\) 上的形式簇范畴 \(\operatorname{Vaf}/k\) 中的群对象。设 \(A\) 为 \(G\) 的仿射代数。\(G\) 的合成律显然定义对角态射，即预有限 \(k\)-代数同态 \[
  \Delta_A:A\longrightarrow A\widehat\otimes_kA.
\]。此同态具有以下性质。

  % SGA3-VII-DIAG-099
\begin{enumerate}\item[\textup{\SGAInlineRefTarget{sga3:VIIB:2.1:item:i}(i)}]图\begin{center}\phantomsection \label{sga3:VIIB:2.1:coassociativity-square} % SGA3-VII-DIAG-099
% Authority: Exp7B-13oct24.pdf, local p. 43.
% Crop: figures/Exp7B_localp043_coassociativity_square.png
% Locator: tex/components/36_expose_VIIB_section21_formal_group_cogroups_en.tex, Loop 1 line 36.
\begin{tikzcd}[row sep=3em,column sep=4.7em]
A \arrow[r,"\Delta_A"] \arrow[d,"\Delta_A"']
  & A\widehat\otimes_k A
      \arrow[d,"{\Delta_A\widehat\otimes\operatorname{id}_A}"] \\
A\widehat\otimes_k A
  \arrow[r,"{\operatorname{id}_A\widehat\otimes\Delta_A}"]
  & A\widehat\otimes_k A\widehat\otimes_k A
\end{tikzcd}
\end{center}交换。

\item[\textup{(ii)}]存在必然唯一的\emph{增广}，即预有限 \(k\)-代数同态 \[
  \varepsilon_A:A\longrightarrow k,
  % SGA3-VII-DIAG-100
\]，使得下面显示的两个复合均为 \(A\) 的恒等映射：\begin{center}\phantomsection \label{sga3:VIIB:2.1:counit-compositions} % SGA3-VII-DIAG-100
% Authority: Exp7B-13oct24.pdf, local p. 43.
% Crop: figures/Exp7B_localp043_counit_compositions.png
% Locator: tex/components/36_expose_VIIB_section21_formal_group_cogroups_en.tex, Loop 1 line 52.
\(
\begin{gathered}
A\xrightarrow{\ \Delta_A\ }A\widehat\otimes_k A
 \xrightarrow{\ \varepsilon_A\widehat\otimes\operatorname{id}_A\ }
 k\widehat\otimes_k A\simeq A
\\[.7ex]
\text{及}\qquad
A\xrightarrow{\ \Delta_A\ }A\widehat\otimes_k A
 \xrightarrow{\ \operatorname{id}_A\widehat\otimes\varepsilon_A\ }
 A\widehat\otimes_k k\simeq A
\end{gathered}
\)
\end{center}

\item[\textup{\SGAInlineRefTarget{sga3:VIIB:2.1:item:iii}(iii)}]存在必然唯一的\emph{对极}，即预有限 \(k\)-代数同态 \[
  c_A:A\longrightarrow A,
  % SGA3-VII-DIAG-101
\]，使得复合\begin{center}\phantomsection \label{sga3:VIIB:2.1:antipode-composition} % SGA3-VII-DIAG-101
% Authority: Exp7B-13oct24.pdf, local p. 44.
% Crop: figures/Exp7B_localp044_antipode_composition.png
% Locator: tex/components/36_expose_VIIB_section21_formal_group_cogroups_en.tex, Loop 1 line 65.
\(
A\xrightarrow{\ \Delta_A\ }A\widehat\otimes_k A
\xrightarrow{\ c_A\widehat\otimes\operatorname{id}_A\ }
A\widehat\otimes_k A
\xrightarrow{\ m_A\ }A
\)
\end{center}等于 \(\eta_A\circ\varepsilon_A\)。这里 \(m_A\) 是将 \(a\widehat\otimes b\) 映到 \(ab\) 的线性映射，而 \(\eta_A:k\to A\) 将 \(\lambda\) 映到 \(\lambda1_A\)。\end{enumerate}

反之，满足 \textup{\SGARefLink{sga3:VIIB:2.1:item:i}{(i)}}--\textup{\SGARefLink{sga3:VIIB:2.1:item:iii}{(iii)}} 的数据 \[
  (\Delta_A,\varepsilon_A,c_A)
\] 赋予 \(G\) 形式 \(k\)-群结构。\footnote{\label{sga3:VIIB:nde:76}也称 \(A\) 为预有限 \(k\)-代数范畴中的\emph{余群}。编者扩充了下文内容，特别是明确了形式群态射 \(K\to G\) 的含义；参见 \SGARefLink{sga3:VIIB:2.3.1}{命题~2.3.1}。}具体地，对每个预有限 \(k\)-代数 \(B\)，连续 \(k\)-模态射 \(\phi:A\to B\) 的集合 \[
  \operatorname{Hom}_c(A,B)
\] 具有关于 \(B\) 函子性的群结构，由 \[
  \phi\cdot\phi'
  =
  m_B\circ
  (\phi\widehat\otimes\phi')
  \circ\Delta_A.
\] 定义。单位元为 \(\eta_B\circ\varepsilon_A\)，其中 \(m_B\) 是 \(B\) 的乘法，\(\eta_B:k\to B\) 将 \(\lambda\) 映到 \(\lambda1_B\)；\(\phi\) 的逆元为 \(\phi\circ c_A\)。连续 \(k\)-代数态射 \(A\to B\) 的集合 \[
  \operatorname{Hom}_{\operatorname{Alp}/k}(A,B)
\] 是一个子群，因为 \(B\) 交换。

\medskip\noindent\textbf{定义。}\phantomsection \label{sga3:VIIB:2.1.definition}\emph{形式 \(k\)-群态射} \(\theta:K\to G\)，按定义就是保持群结构的形式 \(k\)-簇态射。设 \(B\)、\(A\) 分别为 \(K\)、\(G\) 的仿射代数，\(f:A\to B\) 为对应于 \(\theta\) 的态射。此条件等价于 \(f\) 与余乘法相容：\[
  (f\widehat\otimes f)\circ\Delta_A
  =
  \Delta_B\circ f.
\]。此时等式 \[
  \varepsilon_B\circ f=\varepsilon_A,
  \qquad
  c_B\circ f=f\circ c_A
\] 自动成立。以 \(\operatorname{Grf}/k\) 表示形式 \(k\)-群范畴。

\medskip\noindent\textbf{记号。}\phantomsection \label{sga3:VIIB:2.1.notation}下文称理想 \[
  I_A=\operatorname{Ker}(\varepsilon_A)
\] 为 \(A\) 的\emph{增广理想}，并以 \[
  \omega_{G/k}=I_A/\overline{I_A^2}
\] 表示 \(I_A\) 对所有乘积 \(xy\)（其中 \(x,y\in I_A\)）生成的闭理想取商所得的伪紧 \(k\)-模。
